\documentclass[11pt,a4paper]{coverletter}

\senderblock
  {Radheem Bin Razi}
  {Distributed Systems \& Cloud-Native Engineer}
  {Ilmenau, Germany \textbar\ \href{mailto:sheikh.radheem@gmail.com}{sheikh.radheem@gmail.com} \textbar\ \href{https://radheem.github.io/my-cv}{portfolio}}

\begin{document}

%% ===== ENGLISH =====
\recipient{Fraunhofer Institute for Production Technology (IPT)}{Hiring Team}{}

\opening{Dear Hiring Team,}

Bridging physics-based simulation with data-driven prediction lets engineers test process stability without waiting for hours of compute time, and accelerating that cycle is exactly why I want to join the High-Performance Cutting department. I am applying for the Master Thesis position on Machine Learning-Based Methods as Surrogate for Finite Element Modelling.

My recent work has focused on building reliable data pipelines that feed machine learning models with clean, structured inputs. I designed a metrics ingestion system for a 5G testbed where Python applications published per-device measurements to a Kafka pipeline. That pipeline fanned each message into three separate storage targets, including InfluxDB and a graph-compatible database, for downstream analysis. That experience taught me how to prepare datasets for model training while preserving the relationships between connected components. I also built reproducible generation workflows where every model run records its exact inputs, seeds, and quality benchmarks. This ensures that experimental results can be audited and repeated without guesswork.

Translating that background to your thesis work means I can hit the ground running on dataset preparation, model implementation, and result validation. I am comfortable writing Python code to prototype algorithms, and I have already structured experiments to compare baseline methods against newer approaches before committing to full-scale training. The documentation and version-controlled tracking I use for every project will carry directly into your research workflow. This keeps the investigation transparent and the final report production-ready.

I am available to start full-time immediately and can also adjust to a working-student schedule if preferred. I currently live in Germany and can relocate to Aachen within two to three weeks. A Blue Card is straightforward to process since I only need an employment contract to initiate the paperwork. Thank you for considering my application, and I look forward to discussing how I can contribute to your simulation acceleration research.

\closing{Sincerely,}{Radheem Bin Razi}

\clearpage
\selectlanguage{ngerman}

%% ===== DEUTSCH =====
\recipient{Fraunhofer Institute for Production Technology (IPT)}{Personalabteilung}{}

\opening{Sehr geehrtes Hiring-Team,}

Die Verbindung physikbasierter Simulationen mit datengetriebenen Prognosen ermöglicht es Ingenieurinnen und Ingenieuren, die Prozessstabilität zu validieren, ohne stundenlange Rechenzeit in Kauf nehmen zu müssen. Genau diese Beschleunigung des Zyklus ist der Grund, warum ich das High-Performance Cutting department unterstützen möchte. Ich bewerbe mich für die Master Thesis position on Machine Learning-Based Methods as Surrogate for Finite Element Modelling.

In meiner jüngsten Arbeit lag der Fokus auf dem Aufbau zuverlässiger Datenpipelines, die Machine-Learning-Modelle mit sauberen, strukturierten Eingabedaten versorgen. Ich entwickelte ein System zur Metriken-Erfassung für ein 5G testbed, bei dem Python-Anwendungen gerätespezifische Messwerte an eine Kafka-Pipeline übertrugen. Diese Pipeline leitete jede Nachricht auf drei separate Speicherziele, darunter InfluxDB und eine graphenkompatible Datenbank, für die nachgelagerte Analyse. Diese Erfahrung hat mir gezeigt, wie man Datensätze für das Modelltraining aufbereitet, ohne die Zusammenhänge zwischen verbundenen Komponenten zu verlieren. Zudem habe ich reproduzierbare Workflows zur Generierung erstellt, bei denen jeder Modelllauf seine exakten Eingaben, Seeds und Qualitätsbenchmarks protokolliert. Dies stellt sicher, dass experimentelle Ergebnisse lückenlos nachvollzogen und ohne Vermutungen wiederholt werden können.

Die Übertragung dieses Hintergrunds auf Ihre Forschungsarbeit bedeutet, dass ich sofort produktiv in die Datenaufbereitung, die Modellimplementierung und die Validierung der Ergebnisse einsteigen kann. Ich bewege mich sicher beim Schreiben von Python-Code zum Prototyping von Algorithmen und habe Experimente bereits so strukturiert, dass Baseline-Methoden vor dem Einstieg in das Training im großen Maßstab mit neueren Ansätzen verglichen werden können. Die Dokumentation und die versionskontrollierte Nachverfolgung, die ich für jedes Projekt nutze, fließen nahtlos in Ihren Forschungsworkflow ein. Dies gewährleistet Transparenz in der Untersuchung und macht den Abschlussbericht sofort veröffentlichungsreif.

Ich stehe ab sofort vollzeitlich zur Verfügung und passe mich bei Bedarf auch gerne einem working-student schedule an. Ich wohne derzeit in Germany und kann innerhalb von zwei bis drei Wochen nach Aachen ziehen. Die Beantragung einer Blue Card ist unkompliziert, da lediglich ein Arbeitsvertrag zur Einleitung des Verfahrens erforderlich ist. Vielen Dank für die Prüfung meiner Bewerbung. Ich freue mich darauf, in einem persönlichen Gespräch zu besprechen, wie ich Ihren Forschungsarbeiten zur Beschleunigung von Simulationen beitragen kann.

\closing{Mit freundlichen Grüßen,}{Radheem Bin Razi}

\end{document}
